\chapter{连续时间傅里叶变换}
\label{chap:ctft}

傅里叶级数把周期信号分解为离散谐波。令周期趋于无穷大，基频随之趋于零，离散谱线过渡为连续频谱，由此得到连续时间傅里叶变换（continuous-time Fourier transform, CTFT）。以下先给出定义和基本变换对。

\section{从傅里叶级数到傅里叶变换}
\label{sec:ctft-definition}

周期信号的复指数级数为
\begin{equation}
  x(t)=\sum_{k=-\infty}^{\infty}a_k\e^{\iu k\omega_0t},\qquad
  a_k=\frac{1}{T_0}\int_{t_0}^{t_0+T_0}x(t)\e^{-\iu k\omega_0t}\dd t.
  \label{eq:ctft-periodic-series-start}
\end{equation}
令 $T_0\to\infty$，则 $\omega_0=2\pi/T_0\to0$，频率采样点 $\omega=k\omega_0$ 变得连续。定义
\begin{equation}
  X(\iu\omega)=\int_{-\infty}^{\infty}x(t)\e^{-\iu\omega t}\dd t,
  \label{eq:ctft-definition}
\end{equation}
逆变换为
\begin{equation}
  x(t)=\frac{1}{2\pi}\int_{-\infty}^{\infty}X(\iu\omega)\e^{\iu\omega t}\dd\omega.
  \label{eq:ctft-inverse}
\end{equation}
记号 $x(t)\xleftrightarrow{\ \mathcal F\ }X(\iu\omega)$ 表示一对正、逆傅里叶变换。$\iu\omega$ 是频域变量的记号，画频谱时横轴仍是实频率 $\omega$。

\begin{pitfall}[混淆级数系数与变换]
  傅里叶级数系数 $a_k$ 含有 $1/T_0$，只在离散谐波 $k\omega_0$ 上定义；傅里叶变换 $X(\iu\omega)$ 是连续频率函数。过渡时必须同时记住 $a_k=X(\iu k\omega_0)/T_0$ 和求和到积分的替换。
\end{pitfall}

\section{基本变换对与典型计算}
\label{sec:ctft-basic-pairs}

以下变换对直接由定义、逆变换或分布意义下的极限得到。$\Sa(u)=\sin(u)/u$ 沿用第二章的非归一化定义。

\begin{example}[右侧指数衰减]
  对 $a>0$，令 $x(t)=\e^{-at}\stepfun(t)$，则
  \[
    X(\iu\omega)=\int_0^{\infty}\e^{-(a+\iu\omega)t}\dd t
    =\frac{1}{a+\iu\omega}
    =\frac{a-\iu\omega}{a^2+\omega^2}.
  \]
  因而 $|X(\iu\omega)|=1/\sqrt{a^2+\omega^2}$，这是低通型频谱的基本例子。
\end{example}

\begin{example}[单位冲激与直流常数]
  冲激抽样性质给出
  \begin{equation}
    \impulse(t)\xleftrightarrow{\ \mathcal F\ }1,
    \qquad
    1\xleftrightarrow{\ \mathcal F\ }2\pi\impulse(\omega).
    \label{eq:ctft-impulse-constant-pairs}
  \end{equation}
  第二式说明直流分量集中在 $\omega=0$，频谱要用冲激表示。
\end{example}

\begin{example}[对称矩形脉冲]
  设 $x(t)=E$（$|t|<\tau/2$），其余为零。直接积分得到
  \begin{align*}
    X(\iu\omega)&=E\int_{-\tau/2}^{\tau/2}\e^{-\iu\omega t}\dd t
    =\frac{2E\sin(\omega\tau/2)}{\omega}
    =E\tau\Sa\left(\frac{\omega\tau}{2}\right).
  \end{align*}
  因而
  \begin{equation}
    E\,\mathbf{1}_{|t|<\tau/2}
    \xleftrightarrow{\ \mathcal F\ }
    E\tau\Sa\left(\frac{\omega\tau}{2}\right).
    \label{eq:ctft-rectangle-pair}
  \end{equation}
  脉冲越窄，频谱主瓣越宽；脉冲面积 $E\tau$ 等于 $X(0)$。
\end{example}

对应的逆向变换对为
\begin{equation}
  \frac{\sin(\omega_c t)}{\pi t}
  \xleftrightarrow{\ \mathcal F\ }
  \begin{cases}1,&|\omega|<\omega_c,\\0,&|\omega|>\omega_c,\end{cases}
  \label{eq:ctft-sinc-rectangle-pair}
\end{equation}
边界点的取值不影响普通积分意义下的变换。

\begin{example}[单位阶跃的变换]
  单位阶跃不满足绝对可积条件。用截断积分并在分布意义下取极限，得到
  \begin{equation}
    \stepfun(t)\xleftrightarrow{\ \mathcal F\ }
    \frac{1}{\iu\omega}+\pi\impulse(\omega).
    \label{eq:ctft-step-pair}
  \end{equation}
  其中 $1/(\iu\omega)$ 按主值理解，$\pi\impulse(\omega)$ 对应阶跃的直流成分。
\end{example}

\begin{example}[有限区间矩形的尺度计算]
  若矩形脉冲宽度为 $6$、高度为 $2$，由\cref{eq:ctft-rectangle-pair} 得
  $X(\iu\omega)=4\sin(3\omega)/\omega=12\Sa(3\omega)$；若宽度为 $2$、高度为 $1$，则 $X(\iu\omega)=2\Sa(\omega)$。先确定宽度和面积，再代入矩形变换对，比重新积分更不易出错。
\end{example}

线性、时移、调制、微分和卷积性质将在下一节推导。

\section{傅里叶变换的基本性质}
\label{sec:ctft-properties-before-integration}

本节沿用\cref{eq:ctft-definition,eq:ctft-inverse} 的变换约定。普通积分不收敛时，含冲激的变换对仍按广义函数理解。

\subsection{线性性质}

若 $x_1(t)\fourierpair X_1(\iu\omega)$、$x_2(t)\fourierpair X_2(\iu\omega)$，则
\begin{equation}
  a x_1(t)+b x_2(t)
  \fourierpair
  aX_1(\iu\omega)+bX_2(\iu\omega).
  \label{eq:ctft-linearity}
\end{equation}
证明只需把线性组合代入傅里叶积分并拆开积分。计算复杂波形时，通常先将其写成已知变换对的缩放与叠加，再分别变换。

\begin{example}[符号函数的傅里叶变换]
  由 $\operatorname{sgn}(t)=2\stepfun(t)-1$，结合\cref{eq:ctft-step-pair,eq:ctft-impulse-constant-pairs}，得到
  \[
    \operatorname{sgn}(t)
    \fourierpair
    \frac{2}{\iu\omega}.
  \]
  两个冲激项正好抵消。处理阶跃与常数的线性组合时，应把零频冲激一并带入，不能只计算 $1/(\iu\omega)$ 部分。
\end{example}

\subsection{时移性质}

若 $x(t)\fourierpair X(\iu\omega)$，则
\begin{equation}
  x(t-t_0)
  \fourierpair
  X(\iu\omega)\e^{-\iu\omega t_0}.
  \label{eq:ctft-time-shift}
\end{equation}
令 $t'=t-t_0$，直接代入定义可得
\[
  \int_{-\infty}^{\infty}x(t-t_0)\e^{-\iu\omega t}\dd t
  =\e^{-\iu\omega t_0}X(\iu\omega).
\]
乘子 $\e^{-\iu\omega t_0}$ 的模恒为 $1$，所以时移\red{不改变幅度谱}，只给相位谱增加线性项 $-\omega t_0$。若频谱在某些频率处为零，相位在这些点没有定义。

\begin{example}[延迟的矩形脉冲]
  宽度为 $2$、中心位于 $t=1$ 的单位矩形可写成 $\mathbf 1_{0<t<2}$。由\cref{eq:ctft-rectangle-pair,eq:ctft-time-shift}，
  \[
    \mathbf 1_{0<t<2}
    \fourierpair
    2\Sa(\omega)\e^{-\iu\omega}.
  \]
  幅度谱与中心在原点的矩形相同，相位中多出 $-\omega$。
\end{example}

\subsection{频移与调制性质}

若 $x(t)\fourierpair X(\iu\omega)$，则
\begin{equation}
  x(t)\e^{\iu\omega_0t}
  \fourierpair
  X\bigl(\iu(\omega-\omega_0)\bigr).
  \label{eq:ctft-frequency-shift}
\end{equation}
时域乘以复指数，会把整个频谱向右平移 $\omega_0$。利用余弦的复指数展开，还可得到
\begin{equation}
  x(t)\cos(\omega_0t)
  \fourierpair
  \frac12X\bigl(\iu(\omega-\omega_0)\bigr)
  +\frac12X\bigl(\iu(\omega+\omega_0)\bigr).
  \label{eq:ctft-cosine-modulation}
\end{equation}
余弦调制产生位于 $\pm\omega_0$ 的两份频谱副本，每份幅度减半。画图时应先平移横坐标，再乘系数 $1/2$。

\begin{example}[正弦和余弦的冲激频谱]
  从 $1\fourierpair2\pi\impulse(\omega)$ 出发，应用频移与线性性质可得
  \begin{align*}
    \cos(\omega_0t)
      &\fourierpair
      \pi\bigl[\impulse(\omega-\omega_0)+\impulse(\omega+\omega_0)\bigr],\\
    \sin(\omega_0t)
      &\fourierpair
      \frac{\pi}{\iu}
      \bigl[\impulse(\omega-\omega_0)-\impulse(\omega+\omega_0)\bigr].
  \end{align*}
  冲激位于正、负两个频率。实余弦的两条谱线同号，实正弦的两条谱线符号相反。
\end{example}

\subsection{时域微分与频域微分}

设 $x(t)$ 分段光滑，并且分部积分产生的边界项为零，则
\begin{equation}
  \frac{\dd x(t)}{\dd t}
  \fourierpair
  \iu\omega X(\iu\omega).
  \label{eq:ctft-time-differentiation}
\end{equation}
连续使用该性质，得到
\begin{equation}
  \frac{\dd^n x(t)}{\dd t^n}
  \fourierpair
  (\iu\omega)^nX(\iu\omega).
  \label{eq:ctft-higher-time-differentiation}
\end{equation}
例如 $\impulse(t)\fourierpair1$，所以 $\impulse^{(n)}(t)\fourierpair(\iu\omega)^n$。时域微分在频域中会放大高频分量，这是理想微分器容易放大高频噪声的原因。

对\cref{eq:ctft-definition} 关于 $\omega$ 求导，可得另一组关系：
\begin{equation}
  t x(t)
  \fourierpair
  \iu\frac{\dd X(\iu\omega)}{\dd\omega}.
  \label{eq:ctft-frequency-differentiation}
\end{equation}
由 $\e^{-at}\stepfun(t)\fourierpair1/(a+\iu\omega)$（$a>0$），立刻得到
\begin{equation}
  t\e^{-at}\stepfun(t)
  \fourierpair
  \frac{1}{(a+\iu\omega)^2}.
  \label{eq:ctft-ramp-exponential-pair}
\end{equation}
继续求导可推广为 $t^n\e^{-at}\stepfun(t)\fourierpair n!/(a+\iu\omega)^{n+1}$。

\begin{pitfall}[混淆两种微分性质]
  时域求导对应频域乘以 $\iu\omega$；时域乘以 $t$ 对应频域求导后再乘 $\iu$。先看运算发生在哪个域，再选公式。
\end{pitfall}

\subsection{卷积性质}

若 $x(t)\fourierpair X(\iu\omega)$、$h(t)\fourierpair H(\iu\omega)$，则
\begin{equation}
  x(t)\conv h(t)
  \fourierpair
  X(\iu\omega)H(\iu\omega).
  \label{eq:ctft-convolution-property}
\end{equation}
证明从卷积积分出发：
\begin{align*}
  Y(\iu\omega)
  &=\int_{-\infty}^{\infty}\int_{-\infty}^{\infty}
    x(\tau)h(t-\tau)\e^{-\iu\omega t}\dd\tau\dd t\\
  &=\left[\int_{-\infty}^{\infty}x(\tau)\e^{-\iu\omega\tau}\dd\tau\right]
    \left[\int_{-\infty}^{\infty}h(t')\e^{-\iu\omega t'}\dd t'\right]\\
  &=X(\iu\omega)H(\iu\omega).
\end{align*}
这一步默认可以交换积分次序；常见的绝对可积信号满足该条件。卷积在时域往往需要分段，变到频域后只剩乘法，因此\red{卷积可转为频谱相乘}。

\begin{example}[矩形卷积与三角脉冲]
  令 $r(t)=\mathbf 1_{|t|<1}$。由\cref{eq:ctft-rectangle-pair}，
  \[
    r(t)\fourierpair2\Sa(\omega).
  \]
  $r\conv r$ 是支撑在 $[-2,2]$ 上、峰值为 $2$ 的三角脉冲。卷积性质给出
  \[
    (r\conv r)(t)\fourierpair4\Sa^2(\omega).
  \]
  若第二个矩形高度改为 $1/2$，时域三角脉冲的峰值和频谱都相应减半。
\end{example}

\section{LTI 系统的频率响应}
\label{sec:ctft-lti-frequency-response}

连续时间 LTI 系统满足 $y=x\conv h$。由\cref{eq:ctft-convolution-property}，
\begin{equation}
  Y(\iu\omega)=X(\iu\omega)H(\iu\omega),
  \label{eq:ctft-lti-input-output}
\end{equation}
其中 $H(\iu\omega)$ 是单位冲激响应 $h(t)$ 的傅里叶变换，称为系统的频率响应（frequency response）。在 $X(\iu\omega)\ne0$ 的频率上，可以由输入和输出反求
\begin{equation}
  H(\iu\omega)=\frac{Y(\iu\omega)}{X(\iu\omega)}.
  \label{eq:ctft-frequency-response-ratio}
\end{equation}
若 $X$ 在某段频率上为零，这一次实验无法确定该频段的 $H$；直接变换 $h(t)$ 更稳妥。

常见系统的频率响应可由已有变换对直接写出：
\begin{align*}
  h(t)=\impulse'(t)
  &\quad\Longrightarrow\quad H(\iu\omega)=\iu\omega,
  &&\text{理想微分器},\\
  h(t)=\stepfun(t)
  &\quad\Longrightarrow\quad H(\iu\omega)=\frac{1}{\iu\omega}+\pi\impulse(\omega),
  &&\text{从负无穷起算的积分器},\\
  h(t)=\impulse(t-t_0)
  &\quad\Longrightarrow\quad H(\iu\omega)=\e^{-\iu\omega t_0},
  &&\text{理想延迟器}.
\end{align*}
理想延迟器的幅度响应恒为 $1$，相位为 $-\omega t_0$。它只移动波形，不改变各频率分量的幅度。
